\section{Evaluation design}
\label{sec:eval}

\paragraph{Released benchmark data.} SWE-Skills-Bench released pass rates with and without each skill for 98 skills, four harness-model configurations and ten requirement batches per skill, \NSweCells{} skill-by-configuration cells in all \citep{han2026sweskillsbench}. We take $D$ as the difference in batch pass rate. \NSweTies\% of the differences are exactly zero and their standard deviation is \NSweSD. SkillsBench released, for each of \NSbTasks{} tasks and 18 configurations, the number of passes in three trials with and without the task's curated skill \citep{li2026skillsbench}; we take $D=(\text{passes with}-\text{passes without})/3$, which has mean \NSbLift{} points, standard deviation \NSbSD{} and \NSbTies\% zeros. The source paper reports a lift of \PCurated{} points and a no-skill pass rate of \PBase\% over its full evaluation; the released registry gives \NSbLift{} points and \NSbPZero\%. The simulation model uses the published figures, and the reanalyses use the registry. The self-generated condition is not in the released registry. We use these data three ways: to apply admission rules to real per-skill evidence (E8, E9), to build a resampled null with the real tie and magnitude structure by randomizing the signs of the SWE-Skills-Bench differences (E8), and to calibrate and check the simulation model.

\paragraph{Simulation model.} Variant $i$ has difficulty $\theta_i\sim N(0,s_\theta^2)$ on the logit scale; the comparator passes with probability $\operatorname{logit}^{-1}(b_0+\theta_i)$ and the skill arm with $\operatorname{logit}^{-1}(b_0+g_i+\theta_i)$, $g_i\sim N(g,s_g^2)$. With probability $\rho$ both arms share one sampling uniform (common random numbers). $b_0$ is set so the comparator pass rate is the SkillsBench no-skill average, \PBase\%, and $g$ is solved by quadrature so that $\Delta$ equals a target exactly. The defaults $s_\theta=\PSdTheta$, $s_g=\PSdLift$ and $\rho=\PRho$ have no published values; we sweep each. At the SkillsBench average lift with $r=3$ the model gives $\operatorname{sd}(D)=\NSbModelSD$ and \NSbModelTies\% zeros, against \NSbSD{} and \NSbTies\% in the released data, so the model is somewhat less noisy than the benchmark and our cost estimates from it are optimistic. To bound the effect, E14 fits $(s_\theta,s_g,\rho)$ at $r=3$ to three within-skill moments of the released differences: the tie share, the standard deviation around each skill's mean, and the share with $|D|\ge2/3$ (\NRvObsTie\%, \NRvObsSD{} and \NRvObsBig\%). The best fit, $s_\theta=\NRvFitSdTheta$, $s_g=\NRvFitSdLift$ and $\rho=\NRvFitRho$, gives \NRvFitTie\%, \NRvFitSD{} and \NRvFitBig\%; it needs far more heterogeneous skill effects and stronger common random numbers than the defaults. E14 also lets a variant's easiness $\theta_i$ and its skill effect $g_i$ correlate, with correlation $\kappa\in\{-0.5,0,0.5\}$, and reruns the cost claims under each model. Section~\ref{sec:res-cost} also resamples the released data directly.

\paragraph{Gates compared.} We compare seven gates. Two are what a careful benchmark would use today: a one-sided paired $t$-test at the fixed $n=\NFixedN$ that the normal approximation $n=\lceil((z_{1-\alpha}+z_{1-\beta})\,\sigma_D/(\Delta_1-\delta))^2\rceil$ gives for $\Delta_1=\PLiftOne$ points at 80\% power (G1), and the same test re-run after every variant from $n=10$ (G2). G3 admits when the lower end of a two-sided 95\% normal interval exceeds $\delta$, checked every ten variants; this is how ACES- or skillcheck-style intervals behave if a registry reads them as a decision rule. The valid alternatives are the discordant-pair SPRT with its alternative set at the design lift (G4), the normal-mixture mSPRT of \citet{johari2022always} (G5), the betting test with the level-tuned grid bet of Theorem~\ref{thm:adaptive} (G7, the deployed gate, with and without a futility margin), and the betting test with the plug-in bet of Eq.~\eqref{eq:wealth} (G6). Unless stated, ``the betting gate'' below is G7. For G6 we vary the cap ($c=\PCap$, $\PCapAlt$, and the simpler $\lambda_i\le m_0$), the bet (predictable plug-in or aGRAPA, \citealp{waudbysmith2024betting}) and the futility margin. Every sequential gate uses $\alpha=\PAlpha$, $\alpha_{\mathrm f}=\PAlphaF$ and $N_{\max}=\PNmax$. Null cells use \PReps{} replicates, which puts the standard error of a 5\% rate at 0.15 points; power cells use \PRepsPower.

\paragraph{Experiments.} E1 measures false admission under nulls; E2 power and expected variants; E3 library-level FDR over \PWindows{} windows of \PWindow{} candidates with \PNullShare\% nulls, and a resubmission attack; E4 the bias of reported lifts; E5 libraries whose lifts follow the published benchmark summaries; E6 leakage checks; E7 credential cost; E8 and E9 the benchmark reanalyses; E20 compares e-LOND with stopped e-BH and e-BH on adjusted suprema on the E3 histories. Every parameter, with its source or the reason for its value, is in the released \texttt{params.yaml}; every number in Sections~\ref{sec:results} and~\ref{sec:discussion} is generated from the result files.
